% Dot batch 72; source SHA-256: 29760c2aa869bc15f2a63e0985f7fb3c86414392967dd8460eebdcee6fc4c0c2
% Archived source: dot_updates_20261124/batches/batch072/EGA第五卷_正交群挠子与局部等平凡性.tex
% Original body lines: 21--48; integration map: integration/late_integration.json
\begingroup
\def\newscope{V-Q}
\def\OO{\mathcal O}
\def\Qsp{\mathcal Q}
\def\Sym{\operatorname{Sym}}
\def\Isom{\operatorname{Isom}}
\def\Spec{\operatorname{Spec}}
\noindent\textbf{作者按语。}\quad 本可一开始就把16.11.2作为16.2或16.3的推论给出（这两条原本应当互换次序）。这样，本节不依赖“校正判别式”的部分就能集中放在开头。

另外，校正判别式的概念以及16.9在多大程度上已经为人所知，你应该比我清楚，好把功劳归给应当归功的人。也许已经有别的通用术语？我还想指出，真正配叫判别式的是\emph{校正}判别式，因为能推广为任意形式的判别式的是它（参见下一节）。至于二次型的“判别式”（采用这里的术语；如果我没弄错，这也是通用术语，不过手头没有布尔巴基的书可供核对），倒应称为二次型的\emph{行列式}，不应叫判别式。很想听听你对此的看法。如果你同意，我们就借这个机会纠正这点容易误导人的通用术语（直至前几天，我还在不知不觉中把判别式和行列式混在一起）。\cmt{L18-072-01}{这些是作者向收信人提出的术语与归功问题，不是已经查定的历史结论。法文是“通用／公认的术语”；英译写作“recent terminology”。一般齐次形式判别式的下节引用按手稿保留，不用现存附录没有该节而删除。}

\egas{16.12}{评注}
令绝对群概形 $\mathrm{GL}(n)=\mathrm{Aut}(\mathbb Z^n)$ 作用于 $\Qsp(\mathbb Z^n)$。标准二次型的稳定子记为 $\mathrm O(n)$，称为“绝对正交群”。更一般地，对任一二次型 $Q\in\Gamma(S,\Sym^2E)$，引入 $\mathrm{GL}(E)$ 中使 $Q$ 保持不变的闭子群预概形，记作 $\mathrm O(Q)$，称为相对于 $Q$ 的正交群概形。若 $Q$ 同构于 $n$ 变量标准型的基变换，则 $\mathrm O(Q)\simeq\mathrm O(n)_S$。\cmt{L18-072-02}{法文标题为 \emph{Remarques}，英译作 Lemma。法文明列闭子群，并以 $Q$ 为对应的二次型，英译此处有脱漏和 $Q/O$ 混写。}

于是，对任一 $S$ 上的 $S'$，态射 $(**)$ 源中的两点 $u,u'$ 在靶上有相同像，当且仅当 $u'=uv$，其中 $v\in\mathrm O(n)(S')$。\cmt{L18-072-03}{法文在两点记号上方手写 $Q,Q'$，但这里两点在同构空间中，且随后使用 $u'=uv$，故正文统一为 $u,u'$。}
结合16.11，并采用将来要在 EGA 第五或第六卷中更详细论述的术语（眼下可参见 SGAD 第四讲），\cmt{L18-072-04}{法文页边补有“ou en attendant SGAD IV”；英译在此处保留编辑者的待问标记。}源上的自然右 $\mathrm O(n)_S$ 作用连同投影 $(**)$，使 $\Isom(E_0,E)$ 成为 $\Qsp(E)^*$ 上结构群为 $\mathrm O(n)_{\Qsp(E)^*}$ 的主齐性纤维丛。

特别地，$\mathrm{GL}(n)$ 是 $\Qsp(n)^*=\Qsp(\mathbb Z^n)^*$ 上结构群为 $\mathrm O(n)_{\Qsp(n)^*}$ 的主齐性纤维丛。于是（仍取16.9证明中的记号），$P$ 实际上是 $\mathrm O(n)_S$ 下的主齐性纤维丛，且典范地与 $Q$ 相联系：对二次型的同构具有函子性，并与基变换相容。再把纯形式论证与 EGA 的平坦下降理论相结合，就可证明函子 $(E,Q)\mapsto P$ 给出两纤维范畴之间的等价。一边是任意预概形 $S$ 上秩为 $n$ 的局部自由模上的光滑二次型，另一边是这些 $S$ 上结构群为 $\mathrm O(n)_S$ 的主齐性纤维丛。\cmt{L18-072-05}{原参考位置有手改；英译另记需向作者核对。正文保留“EGA 的平坦下降理论”，不替手稿补定卷次。}

\egap{16.13}
设 $Q$ 是有限正秩局部自由模 $E$ 上的光滑二次型。对 $s\in S$，除 $k(s)$ 的特征为 $2$ 且 $E$ 在 $s$ 处的秩为奇数这一情形外，$\mathrm O(Q)$ 在 $s$ 的某个开邻域上光滑；在该例外情形下则不光滑。若这一情形在 $S$ 的所有点都不发生，则（并且仅当如此）上文所考虑的 $P$ 在 $S$ 上光滑：它是 $\mathrm O(n)_S$ 下的主齐性纤维丛，故其光滑性等价于结构群的光滑性。\cmt{L18-072-06}{讨论 $P$ 和条件(iii)时仍用16.9的恒秩 $n$ 约定；秩不恒定时，在各秩层分别陈述。}
利用“亨泽尔引理”，可在16.9的等价条件中加上：
\begin{quote}
\textup{(iii ter)} 同(iii)，但要求 $S'\to S$ 为满射且平展。
\end{quote}

保持此前的假设，还可得到更强的结论。由约化群概形及其主齐性纤维丛的一般理论（参见 SGAD 第二十四讲），若 $Q$ 光滑，则每个 $s$ 都有开邻域 $U$，以及有限、平展、满射的态射 $S'\to U$，使得 $Q_{S'}$ 与标准型同构。比如，若 $S$ 是局部概形，则在(iii ter)中可再要求 $S'\to S$ 有限。

\egap{16.14}
若去掉上述秩与特征的附加假设，16.13的结果便不再普遍成立，首先失效的就是光滑 $Q$ 满足(iii ter)。例如，设 $S=\Spec k$，$k$ 是特征 $2$ 的非完美域，$E$ 的秩为 $1$。取 $a\in k\setminus k^2$，二次型 $at^2$ 不可能在 $k$ 的可分扩张后化为标准型，因为 $k(\sqrt a)/k$ 不可分。\cmt{L18-072-07}{法文明确写 $at^2$，并在“域”上方手加“非完美”；英译在型的式子处讹作 $k^2$。}

不过，一般情形仍可找到有限、局部自由、满射的态射 $S'\to S$，事实上可取 $P/\mathrm{SO}(n)_S$，它是 $\mathrm O(n)_S/\mathrm{SO}(n)_S$ 下的主齐性纤维丛。奇秩时，此商是 $S$ 上的单位平方根群概形 $\mu_{2,S}$；偶秩时，此商是常值群概形 $(\mathbb Z/2\mathbb Z)_S$。\cmt{L18-072-08}{法文手稿将 $\mu_{2,S}\simeq\mathrm O(n)_S/\mathrm{SO}(n)_S$ 补写在“一般情形”中，英译亦未区分奇偶。若按现代定义以 Dickson 不变量的核为偶秩的 $\mathrm{SO}$，必须区分奇偶；参见 Conrad，C.2.11--C.2.12。若本句只承接16.13未涵盖的奇秩情形，原式在该范围内成立。}
基变换 $S'\to S$ 的作用是把结构群从 $\mathrm O(n)$ 降为光滑的 $\mathrm{SO}(n)$，于是前述局部等平凡结论对 $Q_{S'}$ 成立。特别地，若 $S$ 是局部概形，可在(iii bis)中额外要求平凡化标准型的态射有限。
\par\endgroup
